% !TEX root = ../../main.tex

\section{Limitaciones}

La primera limitación es semántica: \textsf{GHC} no prueba que una mónada sea conmutativa. El mecanismo implementado se basa en una declaración explícita del programador mediante \texttt{CD.do}. Si esa declaración es incorrecta, el compilador puede aceptar reordenamientos que cambien efectos observables, como ilustra el control con \texttt{IO}.

La segunda limitación es combinatoria. Enumerar ordenamientos topológicos puede ser costoso cuando el bloque tiene muchos statements independientes. El prototipo permite estudiar exhaustivamente casos pequeños y medianos, pero requiere estrategias adicionales para escalar a programas de mayor tamaño.

La tercera limitación está en la validación. Comparar salidas observables entrega evidencia práctica para los casos evaluados, pero no equivale a una prueba formal de preservación semántica para todos los programas. En el corpus probabilístico, la normalización de distribuciones reduce diferencias representacionales, pero también fija una convención específica de comparación.

Finalmente, el modelo de costos es uniforme por statement. Esta simplificación permite comparar estructura de planes de manera reproducible, pero no distingue entre statements baratos y costosos. Por lo tanto, el menor costo abstracto no necesariamente equivale al menor tiempo de ejecución real.
